% QMB_n06_nogo.tex — Chapter 6: No-Go Theorems
% Sources: Doc. 074, 055
\chapter{No-Go Theorems and What They Show}
\label{ch:nogo}

No-go theorems limit what alternative interpretations of QM
can achieve. The FFGFT must be measured against them.

\section{Bell's Theorem (Revisited)}

Bell's theorem (Chapter~\ref{ch:bell}) excludes local-realistic
theories with hidden variables. The FFGFT circumvents this
not through a loophole, but through a different ontology: The
"hidden variable" is the complete carrier configuration
$(z,r,\theta)_i$ of all modes — but the carrier is shared, not
locally separated.

\section{Kochen-Specker Theorem}

The Kochen-Specker theorem (1967) shows: No definite values can be assigned to a quantum system for all observables simultaneously and
context-independently.\status{E} Measurement results are
context-dependent.

In the FFGFT: The pole transition during measurement depends on the
complete carrier-apparatus interaction, not only on
the mode itself. Context-dependence is thus a property
of the measurement dynamics on the carrier, not a proof of incompleteness
of the description.\status{S} The Born rule is reproduced;
whether the context-dependence is fully derivable from the geometry
remains open.

\section{PBR Theorem}

The Pusey-Barrett-Rudolph theorem (2012) shows: If the quantum
state were only epistemic (expression of knowledge, not of
reality), certain measurements would allow results that
QM excludes.\status{E}

In the FFGFT, the state $(z,r,\theta)$ is a real carrier
configuration, not a state of knowledge — the PBR theorem is thus
compatible: The FFGFT has an ontic wave function.\status{B}

\section{Hardy's Theorem and GHZ State}

Hardy (1992) and Greenberger-Horne-Zeilinger (1989) construct
scenarios in which every local-realistic assignment leads to a contradiction — even without statistical inequalities.\status{E}

Both arguments presuppose separate carriers. In the FFGFT
this presupposition is dropped: The modes share a carrier, and
the GHZ correlations are topologically encoded.\status{S}

\section{What the FFGFT Does Not Claim}

The no-go theorems show what is not possible. The FFGFT claims:
\begin{itemize}
  \item No new predictions in standard Bell tests (all
        statistical results as in QM).\status{B}
  \item A small $\xi$ deviation in high-precision tests
        ($\Delta\mathrm{CHSH}_{\mathrm{elem}}\approx2\times10^{-5}$),
        not measurable today.\status{S}
  \item No violation of freedom of choice (no superdeterminism
        — this approach from Doc.~074 is superseded).
  \item No information faster than light.
\end{itemize}

\quelle{Doc.~055, 074, 230}